\chapter{Arbitrage de latence et sa défense}\label{ch:hf:latency-arbitrage-and-its-defence}

Une cotation qui était juste il y a une milliseconde reste affichée sur une plateforme alors que le prix a déjà bougé sur une autre. La première firme à l'atteindre empoche un
tick sans risque, et le perdant de la course est le teneur de marché qui l'a laissée en place. À la London Stock Exchange, Aquilina, Budish et O'Neill ont compté
de telles courses environ une fois par minute pour chaque valeur du FTSE 100. Le gagnant devançait le premier perdant de cinq à dix microsecondes, et les courses ponctionnaient le trading à hauteur de
0.42 points de base. Dans le modèle de course de ce chapitre, un teneur de marché dont l'annulation prend en médiane 50 microsecondes, face à trois snipers à 40 microsecondes, perd 90~\% de
ses courses. Un retard de 350 microsecondes sur chaque message ne change rien~; le même retard appliqué aux seuls preneurs de liquidité supprime la course. Sur le marché simulé, un
teneur de marché qui lit l'instrument qui mène en moins d'une milliseconde divise par deux la part informée de ses exécutions. Lire plus lentement que l'avance
elle-même ne le protège en rien.

\section{Cotations périmées et courses}

\begin{definition}[Arbitrage de latence, sniping de cotations périmées, course à la latence]\label{def:hf:latency-arbitrage-and-its-defence:arb}
L'\emph{arbitrage de latence}\index{arbitrage de latence} (latency arbitrage) consiste à trader sur un prix que l'information publique a déjà rendu faux, avant que le trader qui l'a affiché
ne puisse le corriger. \emph{sniping de cotations périmées}\index{sniping de cotations périmées} (stale-quote sniping) est sa forme la plus courante~: il s'agit de prendre une cotation affichée qu'un mouvement ailleurs (sur une autre plateforme,
un instrument lié, un communiqué de presse) a laissée périmée. Une \emph{course à la latence}\index{course à la latence} (latency race) est la compétition qui suit chacun de ces événements~:
les ordres des snipers et l'annulation du fournisseur de liquidité partent tous vers le moteur d'appariement, et le premier arrivé décide si la cotation est
prise.
\end{definition}

L'information est publique~; le seul edge est la vitesse. C'est ce qui distingue l'\omterm{def:hf:latency-arbitrage-and-its-defence:arb}{arbitrage de latence} de la sélection adverse du chapitre~6, où le
trader informé sait quelque chose que d'autres ignorent. C'est une conséquence mécanique du traitement des ordres un à un, dans l'ordre de leur arrivée, sur un
carnet d'ordres à cours limité continu (le livre~10, chapitre~1)~: celui qui est plus rapide d'une microseconde rafle tout, et le prix de la course est payé par
la cotation.

\section{Les preuves~: courses, leur ampleur et qui gagne}

Budish, Cramton et Shim ont étudié le contrat à terme S\&P 500 et son fonds de 2005 à 2011. La durée de vie médiane d'une opportunité d'arbitrage entre les deux est passée de
97 millisecondes à 7, tandis que le profit médian par opportunité restait voisin de 0.08 point d'indice~: les investissements dans la vitesse ont raccourci les courses sans réduire
l'enjeu, et celui-ci est revenu au plus rapide. Aquilina, Budish et O'Neill ont utilisé les données de messages de la London Stock Exchange, qui contiennent les ordres et
les annulations échouées aussi bien que celles qui ont abouti, et montrent donc les perdants de chaque course. Sur les valeurs du FTSE 100, ils ont relevé environ 537 courses par jour et par
valeur~; un écart modal entre le gagnant et le premier perdant de 5 à 10 microsecondes~; environ 22~\% du volume échangé au sein des courses~; et une taxe d'\omterm{def:hf:latency-arbitrage-and-its-defence:arb}{arbitrage de latence} de
0.42 points de base de l'ensemble du volume échangé (0.53 hors des courses), face à une fourchette effective légèrement supérieure à 3. Supprimer les courses réduirait
le coût de la liquidité d'environ 17~\%, soit quelque 5~milliards de dollars par an sur les actions mondiales. Shkilko et Sokolov ont mis en évidence l'autre versant~: quand la météo a perturbé les
liaisons hertziennes des traders les plus rapides, la sélection adverse et les coûts de trading ont baissé.

Le modèle de course de ce chapitre (\cref{lst:hf:latency-arbitrage-and-its-defence:race}) reproduit le mécanisme. Après chaque événement, l'annulation du fournisseur de liquidité
et les ordres de trois snipers atteignent le moteur après des latences lognormales~: des médianes de 50 microsecondes pour l'annulation et de 40 pour chaque sniper,
gigue 0.3. La cotation est à un demi-tick et le prix saute d'un tick. Le fournisseur subit un sniping dans 89.9~\% des courses, en perdant
en moyenne 0.45 tick par action et par course. Parmi les snipers, le gagnant devance le premier perdant d'une médiane de 7.3 microsecondes, l'ordre de grandeur mis en évidence à Londres.

\section{La défense par le teneur de marché}

Un teneur de marché dispose de trois défenses. La première est la vitesse~: la \cref{fig:hf:latency-arbitrage-and-its-defence:sniped} montre la probabilité de
subir un sniping selon la latence médiane d'annulation du fournisseur, les snipers étant fixés à 40 microsecondes. À 20 microsecondes, elle vaut 12~\%~; à parité, 75~\%~; à
100, elle est certaine. La deuxième est la prédiction~: retirer une cotation avant l'arrivée de l'information, à partir d'un signal qui bouge en premier (le
skew du chapitre~7, l'avance du chapitre~8). La troisième est la largeur et la taille~: une cotation plus large ou plus petite perd moins à chaque course. En pratique, les trois se combinent et toutes ont un coût~:
la vitesse est un coût fixe (chapitre~1), la prédiction renonce à du volume (chapitre~6), la largeur renonce à la capture.

\begin{omfigure}
\begin{tikzpicture}
\begin{axis}[width=11cm,height=5.2cm,xlabel={\small latence médiane d'annulation du fournisseur (microsecondes)},ylabel={\small probabilité de subir un sniping},
  xmode=log,log basis x=10,xtick={10,20,30,50,70,100},xticklabels={10,20,30,50,70,100},xmin=9,xmax=110,ymin=0,ymax=1.02,
  tick label style={font=\footnotesize},grid=major,grid style={gray!25},table/col sep=comma,
  legend style={font=\scriptsize,at={(0.5,-0.3)},anchor=north},legend columns=2]
\addplot[omThm,thick,mark=*,mark size=1.3pt] table[x=lp_median,y=continuous]{figdata/hft/09-latency-arbitrage-and-its-defence/sniped.csv};
\addplot[omDef,thick,dashed,mark=square*,mark size=1.3pt] table[x=lp_median,y=asymmetric]{figdata/hft/09-latency-arbitrage-and-its-defence/sniped.csv};
\draw[gray,dashed] (axis cs:40,0) -- (axis cs:40,1.02);
\node[font=\scriptsize,anchor=west,fill=white,inner sep=1pt] at (axis cs:41,0.5) {médiane des snipers};
\legend{continu,{preneurs retardés de 50 microsecondes}}
\end{axis}
\end{tikzpicture}

\omcaption{Probabilité qu'une cotation périmée d'un fournisseur de liquidité soit prise avant l'arrivée de son annulation, en fonction de la latence médiane d'annulation du fournisseur
(échelle logarithmique), avec trois snipers à une médiane de 40 microsecondes (gigue lognormale 0.3), sur un carnet continu et sur un carnet qui retarde les preneurs de 50
microsecondes~; 100\,000 courses par point. Données~: \texttt{hf\_latency.by\_lp\_latency}.}\label{fig:hf:latency-arbitrage-and-its-defence:sniped}
\end{omfigure}

La défense par l'avance peut être testée sur le marché simulé. Un teneur de cotations à un lot du fonds, qui suit le future une demi-seconde plus tard, retire pendant une
seconde le côté de sa cotation que le mouvement du future laisse périmé (\cref{lst:hf:latency-arbitrage-and-its-defence:guard}). Les traders informés du
marché du fonds jouent les snipers~: ils tradent vers le prix efficient du fonds lorsqu'il bouge. En faisant varier le retard avec lequel le teneur de cotations lit le future, sur
quatre sessions de vingt minutes~:

\begin{center}
\small
\begin{tabular}{@{}lrrrrrr@{}}
\toprule
lit le future après & sans garde & 1 ms & 100 ms & 300 ms & 600 ms & 1 s\\
\midrule
actions par session & 7\,550 & 5\,225 & 5\,875 & 6\,750 & 7\,425 & 7\,175\\
part informée des exécutions & 28.8~\% & 13.4~\% & 18.7~\% & 20.0~\% & 24.6~\% & 29.6~\%\\
mark-out à cinq secondes (ticks par action) & 0.174 & 0.402 & 0.374 & 0.330 & 0.192 & 0.099\\
\bottomrule
\end{tabular}
\end{center}

Lu en moins d'une milliseconde, le future divise par deux la part informée et plus que double le mark-out
(la \cref{fig:hf:latency-arbitrage-and-its-defence:defence}). Le bénéfice diminue à mesure que la lecture ralentit et disparaît dès qu'elle est plus lente que l'avance
elle-même~: à 600 millisecondes, la garde réagit après que les traders informés ont agi, et à une seconde elle retire des cotations aux mauvais moments.

\begin{omfigure}
\begin{tikzpicture}
\begin{axis}[width=11cm,height=5.2cm,axis y line*=right,axis x line=none,ylabel={\small mark-out (ticks par action)},xmode=log,log basis x=10,
  xmin=0.7,xmax=1400,ymin=0,ymax=0.5,tick label style={font=\footnotesize},table/col sep=comma]
\addplot[omDef,thick,dashed,mark=square*,mark size=1.3pt] table[x=delay_ms,y=markout]{figdata/hft/09-latency-arbitrage-and-its-defence/defence.csv};
\label{plot:hf:guardmo}
\end{axis}
\begin{axis}[width=11cm,height=5.2cm,axis y line*=left,xlabel={\small retard avec lequel le teneur de cotations lit le future (millisecondes)},
  ylabel={\small part informée des exécutions (\%)},xmode=log,log basis x=10,xmin=0.7,xmax=1400,xtick={1,10,100,1000},xticklabels={1,10,100,1000},
  ymin=0,ymax=35,tick label style={font=\footnotesize},grid=major,grid style={gray!25},table/col sep=comma,
  legend style={font=\scriptsize,at={(0.5,-0.3)},anchor=north},legend columns=2]
\addplot[omThm,thick,mark=*,mark size=1.3pt] table[x=delay_ms,y=informed]{figdata/hft/09-latency-arbitrage-and-its-defence/defence.csv};
\addlegendentry{part informée (gauche)}
\addlegendimage{/pgfplots/refstyle=plot:hf:guardmo}\addlegendentry{mark-out à cinq secondes (droite)}
\draw[gray,dashed] (axis cs:500,0) -- (axis cs:500,35);
\node[font=\scriptsize,anchor=east,fill=white,inner sep=1pt] at (axis cs:480,5) {avance du future};
\end{axis}
\end{tikzpicture}

\omcaption{Un teneur de cotations à un lot du fonds qui retire le côté que le mouvement du future laisse périmé~: part informée et mark-out à cinq secondes de ses exécutions
en fonction du retard avec lequel il lit le future (échelle logarithmique)~; le future mène le fonds d'une demi-seconde~; quatre sessions de vingt minutes. Sans la garde~: 28.8~\%
et 0.174. Données~: \texttt{hf\_latency.defence}.}\label{fig:hf:latency-arbitrage-and-its-defence:defence}
\end{omfigure}

\section{Défense par la plateforme~: ralentisseurs et enchères par lots}

\begin{definition}[Ralentisseur, ralentisseur asymétrique]\label{def:hf:latency-arbitrage-and-its-defence:bump}
Un \emph{ralentisseur}\index{ralentisseur} (speed bump) est un délai fixe que la plateforme ajoute aux messages avant qu'ils n'atteignent son moteur d'appariement. Un \emph{ralentisseur asymétrique
}\index{ralentisseur asymétrique} (asymmetric speed bump) ne retarde que certains messages, typiquement ceux qui prennent de la liquidité, tandis que les annulations et les nouvelles cotations passent sans délai, de sorte que
un fournisseur de liquidité puisse corriger une cotation périmée avant qu'un preneur ne puisse l'atteindre.
\end{definition}

\begin{proposition}[Ce que change un délai]\label{prop:hf:latency-arbitrage-and-its-defence:bump}
Si chaque message entrant est retardé de la même durée $d$, l'ordre d'arrivée au moteur, et donc le vainqueur de chaque course, reste inchangé. Si seuls
les preneurs sont retardés de $d$, une cotation périmée n'est prise que lorsque la latence du premier preneur majorée de $d$ est inférieure à la latence d'annulation du fournisseur, ce qui
ne se réalise dans aucune course dès que $d$ dépasse l'écart entre l'annulation la plus lente et le preneur le plus rapide.
\end{proposition}

\begin{proof}
Les instants d'arrivée $t+x_i$ deviennent $t+x_i+d$~: un décalage commun conserve leur ordre. Si seuls les preneurs sont retardés, la comparaison s'écrit $x_{\text{prise}}+d<x_{\text{annulation}}$.
\end{proof}

Un délai symétrique protège contre autre chose~: un délai appliqué aux données sortantes de la plateforme, comme à ses ordres entrants, empêche un trader de
réagir à ses propres exécutions avant que la plateforme ne transmette l'ordre d'un client. Le modèle de course du chapitre chiffre les différentes conceptions~:

\begin{center}
\small
\setlength{\tabcolsep}{4pt}%
\begin{tabular}{@{}>{\raggedright\arraybackslash}p{6.6cm}>{\centering\arraybackslash}p{3.6cm}>{\centering\arraybackslash}p{3.6cm}@{}}
\toprule
conception de la plateforme & probabilité de subir un sniping & taxe par course (ticks par action)\\
\midrule
continu & 0.899 & 0.449\\
chaque message retardé de 350 microsecondes & 0.899 & 0.449\\
preneurs retardés de 350 microsecondes & 0.000 & 0.000\\
enchère par lots fréquente, 100 microsecondes & 0.211 & 0.106\\
enchère par lots fréquente, une milliseconde & 0.021 & 0.011\\
\bottomrule
\end{tabular}
\end{center}

Une enchère par lots fréquente (le livre~10, chapitre~10) traite ensemble tous les ordres qui arrivent pendant un intervalle, à un prix unique, si bien que la vitesse
à l'intérieur de l'intervalle ne vaut rien~; une cotation périmée n'est prise que si les preneurs arrivent dans un lot antérieur à celui de l'annulation. Budish, Cramton et Shim
l'ont proposée comme réponse de conception de marché à la course. Plus le lot est long, plus la perte est rare.

\begin{dated}{2026-09}{Ralentisseurs sur les actions américaines}\label{dat:hf:latency-arbitrage-and-its-defence:bumps}
La SEC a approuvé IEX comme bourse de valeurs nationale le 17 juin 2016 (Release 34-78101), avec un point de présence et une bobine d'environ 38~milles de
fibre optique ajoutant 350 microsecondes aux messages dans les deux sens. En 2020, elle a approuvé l'ordre D-Limit d'IEX (SR-IEX-2019-15), un ordre à cours limité affiché
qu'IEX réajuste lorsque son Crumbling Quote Indicator, un modèle bâti sur les cotations de huit autres bourses, prévoit un changement imminent du
meilleur prix protégé~: une protection asymétrique mise en œuvre par la plateforme. En février 2020, elle a refusé d'approuver la Liquidity Provider Protection de Cboe EDGA, qui
aurait retardé les ordres exécutables entrants de quatre millisecondes, tout en laissant les ordres au repos être annulés ou modifiés sans délai. Le 11 juin
2026, elle a proposé d'abroger la Rule 611 de la Regulation NMS, l'interdiction du trade-through qui fait de la cotation protégée de chaque plateforme une cible pour
le routage inter-plateformes (le livre~1, chapitre~9).
\end{dated}

\section{Le last look comme défense}

Sur le FX de gré à gré, la défense du fournisseur de liquidité est contractuelle. Sous last look (le livre~2, chapitre~15), un fournisseur qui diffuse un prix
peut, pendant un temps de rétention, le comparer au marché et rejeter une demande d'exécution formulée sur un prix qui a évolué à son détriment. C'est un \omterm{def:hf:latency-arbitrage-and-its-defence:bump}{ralentisseur
asymétrique} contrôlé par une seule des parties, et la raison pour laquelle le FX Global Code fixe des principes sur la manière dont il peut être utilisé. Le chapitre~21 price les diffusions de prix avec et sans last look.

\section{Fiches de stratégie}

\begin{strategyfile}{Sniping de cotations périmées entre plateformes}\label{strat:hf:latency-arbitrage-and-its-defence:snipe}
\sfield{Qui paie, et pourquoi} Les fournisseurs de liquidité dont les cotations sur une plateforme n'ont pas encore rattrapé une transaction ou un changement de cotation sur une autre.
\sfield{Instruments et plateformes} Actions négociées sur de nombreuses plateformes~; le même contrat à terme sur deux plateformes~; des cryptomonnaies sur plusieurs bourses.
\sfield{Signal} Un changement de prix sur la plateforme la plus rapide à se mettre à jour, qui laisse une cotation au repos ailleurs du mauvais côté du nouveau prix.
\sfield{Dimensionnement et exécution} Ordres exécuter-ou-annuler pour la taille de la cotation périmée, envoyés dès que le mouvement est repéré.
\sfield{Coûts} Frais du preneur~; les pertes des courses perdues (ordres qui manquent leur cible, ou qui arrivent après la disparition de la cotation et s'exécutent à de moins bons prix)~; le
coût fixe d'être le plus rapide.
\sfield{Comment elle s'épuise} Seuls les plus rapides gagnent~: Aquilina, Budish et O'Neill ont trouvé des gagnants qui devançaient de 5--10 microsecondes~; les plateformes et les fournisseurs se défendent
(retards asymétriques, annulation sur signal).
\sfield{Horizon, capacité, infrastructure} Microsecondes~; faible par course, importante en agrégé (environ 5~milliards de dollars par an sur les actions mondiales)~; colocation,
micro-ondes et matériel (les livres~13 et~14).
\sfield{Backtest honnêtement} Rejeu au niveau des messages avec la latence de chaque concurrent~; un rejeu sans les messages des perdants ne peut pas montrer les courses.
\sfield{Sources} Aquilina, Budish et O'Neill (2022)~; le modèle de course du chapitre~: 89.9~\% des courses gagnées par les snipers face à un fournisseur
plus lent de 10 microsecondes.
\end{strategyfile}

\begin{strategyfile}{Sniping de cotations d'actions mené par les futures}\label{strat:hf:latency-arbitrage-and-its-defence:futures}
\sfield{Qui paie, et pourquoi} Les teneurs de marché de fonds et d'actions dont les cotations sont en retard sur un mouvement du future sur indice négocié dans une autre ville.
\sfield{Instruments et plateformes} Futures sur indices, ainsi que les fonds et les actions de l'indice.
\sfield{Signal} Le mouvement du future, converti dans les unités de chaque instrument.
\sfield{Dimensionnement et exécution} Prendre d'abord les cotations périmées des instruments les plus exposés.
\sfield{Coûts} La liaison la plus rapide entre les deux villes~; frais du preneur.
\sfield{Comment elle s'épuise} La course aux armements~: 97 millisecondes d'opportunité médiane en 2005, 7 en 2011~; quand le mauvais temps a coupé les liaisons par micro-ondes, la sélection adverse
et les coûts de négociation ont baissé (Shkilko et Sokolov).
\sfield{Horizon, capacité, infrastructure} Millisecondes~; liaisons longue distance par micro-ondes et par fibre (le livre~14, chapitres~13--14).
\sfield{Backtest honnêtement} La latence réelle d'une liaison dans un seul sens et ses interruptions météo, ainsi que la vitesse d'annulation propre aux cibles.
\sfield{Sources} Budish, Cramton et Shim (2015)~; Shkilko et Sokolov (2020).
\end{strategyfile}

\begin{strategyfile}{Fourniture de liquidité derrière un ralentisseur asymétrique}\label{strat:hf:latency-arbitrage-and-its-defence:bump}
\sfield{Qui paie, et pourquoi} Les preneurs qui acceptent le délai sur les cotations de la plateforme~; la plateforme protège le fournisseur des plus rapides.
\sfield{Instruments et plateformes} Plateformes à protections asymétriques, comme un ordre affiché réajusté par la plateforme (le D-Limit d'IEX).
\sfield{Signal} Le \omterm{def:hf:fair-value:fair}{prix juste} propre au fournisseur~; la protection de la plateforme gère les courses.
\sfield{Dimensionnement et exécution} Coter plus serré ou plus gros que sur les plateformes continues, puisque moins de cotations font l'objet d'un sniping.
\sfield{Coûts} Moins d'exécutions~: certains preneurs évitent les plateformes retardées~; le réajustement par la plateforme peut déplacer la cotation alors que le fournisseur ne l'aurait pas fait.
\sfield{Comment elle s'épuise} La réglementation~: la SEC a refusé d'approuver un retard asymétrique de quatre millisecondes sur Cboe EDGA en 2020~; les preneurs contournent la plateforme.
\sfield{Horizon, capacité, infrastructure} Comme pour toute cotation~; la plateforme fournit la protection.
\sfield{Backtest honnêtement} Modéliser qui cesse d'acheminer des ordres vers la plateforme~; le zéro sniping du modèle de course à 350 microsecondes suppose que les preneurs continuent de venir.
\sfield{Sources} Ordonnances de la SEC de 2016 et de 2020 (IEX) et de 2020 (EDGA)~; le modèle de course du chapitre.
\end{strategyfile}

\begin{strategyfile}{Annulation défensive de cotations sur un signal d'avance}\label{strat:hf:latency-arbitrage-and-its-defence:guard}
\sfield{Qui paie, et pourquoi} Les preneurs non informés continuent de trader avec le fournisseur~; les snipers ne trouvent aucune cotation périmée.
\sfield{Instruments et plateformes} Tout instrument suiveur ayant un meneur~: un fonds et son future, une action et son secteur, une plateforme secondaire et la
principale.
\sfield{Signal} Le mouvement du meneur moins celui du suiveur sur une courte fenêtre, lu aussi vite que possible.
\sfield{Dimensionnement et exécution} Retirer le côté périmé pour une courte durée~; revenir en fin de file d'attente.
\sfield{Coûts} Volume abandonné (31~\% ici à 1~ms), messages, places dans la file.
\sfield{Comment elle s'épuise} Quand les snipers lisent le meneur plus vite que le fournisseur~: à 600 millisecondes, au-delà de l'avance, la garde ne protège plus.
\sfield{Horizon, capacité, infrastructure} De la microseconde à la milliseconde~; un flux de données de la plateforme du meneur.
\sfield{Backtest honnêtement} La latence réelle du fournisseur vers les données du meneur et celle des snipers, plutôt qu'une horloge fusionnée.
\sfield{Sources} Ce chapitre~: part informée des exécutions de 28.8~\% sans la garde, 13.4~\% avec une lecture à 1~milliseconde, 24.6~\% à 600 millisecondes.
\end{strategyfile}

\section{Tutoriel~: qui gagne la course}\label{tut:hf:latency-arbitrage-and-its-defence:tut}

\begin{tutorial}
\textbf{Objectif.} Simuler les courses de latence sous chaque conception de plateforme, et tester une défense fondée sur le meneur sur le marché simulé. \textbf{État final~:} les deux
tableaux et les \cref{fig:hf:latency-arbitrage-and-its-defence:sniped,fig:hf:latency-arbitrage-and-its-defence:defence}.
\begin{tutsteps}
\item \textbf{Courses.} \texttt{firm.latrace.race} calcule la latence de chaque message et applique la règle de la plateforme.
\omcode{code/firm/latrace/firm_latrace.py}{37}{58}{Cent mille courses simultanées, sous quatre conceptions de plateforme.}\label{lst:hf:latency-arbitrage-and-its-defence:race}
\item \textbf{Conceptions et latences.} \texttt{hf\_latency.designs()} et \texttt{by\_lp\_latency()}.
\item \textbf{La garde.} \texttt{LeadGuard} lit le future avec son propre délai et retire le côté périmé pendant une seconde.
\omcode{code/firm/latrace/firm_latrace.py}{81}{95}{Retirer le côté que le mouvement du meneur laisse périmé.}\label{lst:hf:latency-arbitrage-and-its-defence:guard}
\item \textbf{Test} avec \texttt{defence()} sur quatre sessions où le fonds suit le future avec un retard d'une demi-seconde.
\end{tutsteps}
\textbf{Pour aller plus loin.} Remplacer \texttt{firm\_tape} par \texttt{firm.exchsim} (le livre~10, chapitre~26) avec deux plateformes et une
matrice de latences, et faire courir de vrais agents snipers~; donner aux snipers un avantage de latence aléatoire qui change chaque jour.
\end{tutorial}

\section{Construction~: le simulateur de courses}\label{bld:hf:latency-arbitrage-and-its-defence:latrace}

\begin{build}
\bfield{Rôle} La probabilité de perdre une course et son coût, sous une conception de plateforme, pour un profil de latence donné~; et une garde de meneur pour tout teneur de cotations.
\bfield{Interface}
\texttt{race(design, n, lp\_median, sniper\_median, snipers, jitter, d, tau, jump, half, seed) -> \{p\_sniped, tax, gap\_median, first\_loser\_gap\}}~; \texttt{LeadGuard(leader\_t, leader\_mid, lead\_delay, move, window, hold, size, limit)}.
\bfield{Règles} Les latences sont tirées par message~; conceptions~: continue, symétrique, asymétrique, par lots (les annulations passent en premier au sein d'un lot). La garde lit
le meneur avec son propre délai, jamais la réalité du simulateur.
\bfield{Tests d'acceptation} \texttt{code/firm/latrace/tests/}~: un délai symétrique laisse inchangé le vainqueur de chaque course~; un délai asymétrique de 350~microsecondes
protège chaque cotation~; un fournisseur plus rapide subit moins de sniping~; un lot d'une milliseconde réduit le sniping de plus de quatre-vingt-dix pour cent~; la taxe égale le
produit de la probabilité par l'edge périmé.
\bfield{Pour aller plus loin} Courses sur \texttt{firm.exchsim} avec ses réglages de délai asymétrique et d'intervalle de lot~; un sniper qui apprend la
latence d'annulation du fournisseur.
\end{build}

\begin{omsources}
\item M.~Aquilina, E.~Budish and P.~O'Neill, ``Quantifying the high-frequency trading arms race'', \emph{Quarterly Journal of Economics} 137(1), 2022.
\item E.~Budish, P.~Cramton and J.~Shim, ``The high-frequency trading arms race: frequent batch auctions as a market design response'',
\emph{Quarterly Journal of Economics} 130(4), 2015.
\item A.~Shkilko and K.~Sokolov, ``Every cloud has a silver lining: fast trading, microwave connectivity, and trading costs'', \emph{Journal of
Finance} 75(6), 2020.
\item SEC, Release 34-78101 (IEX, 2016)~; ordonnance approuvant SR-IEX-2019-15 (D-Limit, 2020)~; ordonnance rejetant SR-CboeEDGA-2019-012 (2020)~; fiche d'information sur les réformes de la Regulation
NMS (2026).
\end{omsources}

\section{Exercices}

\begin{exercise}[$\star$]\label{exo:hf:latency-arbitrage-and-its-defence:1}
Une cotation à un demi-tick est prise après un saut de deux ticks. Calculer le gain du sniper par action.
\end{exercise}

\begin{exercise}[$\star$]\label{exo:hf:latency-arbitrage-and-its-defence:2}
Un symbole compte 537 courses par jour sur une séance de huit heures et demie. Calculer le nombre de courses par minute.
\end{exercise}

\begin{exercise}[$\star$]\label{exo:hf:latency-arbitrage-and-its-defence:3}
L'annulation du fournisseur arrive 60 microsecondes après un événement, et l'ordre du preneur le plus rapide 45 microsecondes après. Calculer le délai, imposé aux seuls preneurs, qui protège la cotation.
\end{exercise}

\begin{exercise}[$\star\star$]\label{exo:hf:latency-arbitrage-and-its-defence:4}
Montrer qu'un délai appliqué de la même façon à chaque message entrant ne modifie aucune course, puis indiquer à quoi reste utile un délai symétrique.
\end{exercise}

\begin{exercise}[$\star\star$]\label{exo:hf:latency-arbitrage-and-its-defence:5}
Expliquer pourquoi une enchère par lots fréquente avec un intervalle d'une milliseconde laisse encore 2~\% des cotations périmées être exécutées dans le modèle du chapitre.
\end{exercise}

\begin{exercise}[$\star\star$]\label{exo:hf:latency-arbitrage-and-its-defence:6}
À partir de la \cref{fig:hf:latency-arbitrage-and-its-defence:defence}, pour quels délais de lecture la garde réduit encore de moitié l'écart entre la
part informée sans garde et celle à une milliseconde~?
\end{exercise}

\begin{exercise}[$\star\star\star$]\label{exo:hf:latency-arbitrage-and-its-defence:7}
\emph{Programmation.} Lancer \texttt{race} avec la conception asymétrique, \texttt{lp\_median=70}, \texttt{d=50} et les autres paramètres du chapitre, puis calculer à quelle fréquence le fournisseur subit du sniping, et de
combien le délai des preneurs doit augmenter pour la ramener sous 1~\%~?
\end{exercise}

\begin{exercise}[$\star\star\star$]\label{exo:hf:latency-arbitrage-and-its-defence:8}
\emph{Trouver la faille.} «~Notre rejeu du flux du carnet d'ordres de la bourse ne montre aucune \omterm{def:hf:latency-arbitrage-and-its-defence:arb}{course à la latence} sur cette action, donc l'\omterm{def:hf:latency-arbitrage-and-its-defence:arb}{arbitrage de latence} ne nous touche pas.~»
\end{exercise}

\section{Problème~: la course que personne ne voit}

\begin{problem}[{Problème du week-end --- la course que personne ne voit}]\label{pb:hf:latency-arbitrage-and-its-defence:1}
Un teneur de marché au sein d'un fonds soupçonne qu'il paie pour la vitesse des autres firmes. Déterminer à combien se monte ce coût et ce qui pourrait y mettre fin.

\medskip\noindent\textbf{Partie~I --- Courses.}
\begin{enumerate}
\item Définir l'\omterm{def:hf:latency-arbitrage-and-its-defence:arb}{arbitrage de latence}, le \omterm{def:hf:latency-arbitrage-and-its-defence:arb}{sniping de cotations périmées} et la \omterm{def:hf:latency-arbitrage-and-its-defence:arb}{course à la latence}, et les distinguer
du trading informé.
\item Que Budish, Cramton et Shim ont-ils établi sur l'arbitrage entre le future et le fonds de 2005 à 2011~?
\item Donner les décomptes d'Aquilina, Budish et O'Neill, les écarts entre courses, la part du volume et la taxe.
\item Pourquoi pouvaient-ils voir des courses qu'un flux de carnet d'ordres ne peut pas montrer~?
\end{enumerate}

\medskip\noindent\textbf{Partie~II --- Le modèle.}
\begin{enumerate}[resume]
\item Décrire le modèle de course et ses paramètres.
\item Donner la probabilité d'être victime de sniping et la taxe sur un carnet continu.
\item Donner la probabilité pour des latences du fournisseur de 20, 40 et 100 microsecondes.
\item Quel est l'écart entre le vainqueur et le premier perdant, et comment se situe-t-il par rapport à Londres~?
\end{enumerate}

\medskip\noindent\textbf{Partie~III --- Conceptions des plateformes.}
\begin{enumerate}[resume]
\item Définir un \omterm{def:hf:latency-arbitrage-and-its-defence:bump}{ralentisseur} et un \omterm{def:hf:latency-arbitrage-and-its-defence:bump}{ralentisseur asymétrique}.
\item Démontrer la \cref{prop:hf:latency-arbitrage-and-its-defence:bump}.
\item Donner la probabilité de sniping sous chaque conception du tableau.
\item Résumer le bilan réglementaire de l'encadré daté.
\item Pourquoi le last look est-il un \omterm{def:hf:latency-arbitrage-and-its-defence:bump}{ralentisseur asymétrique}, et qui le contrôle~?
\end{enumerate}

\medskip\noindent\textbf{Partie~IV --- Le verdict.}
\begin{enumerate}[resume]
\item Énoncer le \emph{résultat nommé}~: la taxe d'\omterm{def:hf:latency-arbitrage-and-its-defence:arb}{arbitrage de latence} rapportée à la demi-fourchette du fournisseur sur un carnet continu, avec un délai symétrique,
avec un délai asymétrique et avec un lot d'une milliseconde.
\item Donner la part informée et le mark-out de la garde de meneur à 1, 300 et 600 millisecondes, et expliquer la régularité observée.
\item Quel est le coût de la garde~?
\item Quelle défense retenir pour un nouveau teneur de marché de fonds qui ne dispose pas des liaisons les plus rapides~?
\item Pourquoi les preneurs pourraient-ils cesser d'acheminer leurs ordres vers une plateforme dotée d'un délai asymétrique~?
\item Qu'est-ce que l'abrogation de la règle de trade-through changerait pour le sniping entre plateformes~?
\item En une phrase~: qui paie l'\omterm{def:hf:latency-arbitrage-and-its-defence:arb}{arbitrage de latence}~?
\end{enumerate}
\end{problem}

\section{Questions d'entretien}

\begin{interviewq}[$\star$ \iqroles{trader}]\label{iq:hf:latency-arbitrage-and-its-defence:1}
Quelle est la différence entre subir une sélection adverse et être victime de sniping~?
\end{interviewq}

\begin{interviewq}[$\star\star$ \iqroles{chercheur}]\label{iq:hf:latency-arbitrage-and-its-defence:2}
Comment mesureriez-vous la fréquence à laquelle vos cotations perdent des \omterm{def:hf:latency-arbitrage-and-its-defence:arb}{courses à la latence}, à partir de vos propres données~?
\end{interviewq}

\begin{interviewq}[$\star\star$ \iqroles{développeur}]\label{iq:hf:latency-arbitrage-and-its-defence:3}
Où chercheriez-vous à gagner dix microsecondes sur votre chemin d'annulation~?
\end{interviewq}

\begin{interviewq}[$\star\star$ \iqroles{chercheur}]\label{iq:hf:latency-arbitrage-and-its-defence:4}
Pourquoi une enchère par lots fréquente supprime-t-elle la valeur de la vitesse au sein du lot, mais pas à ses frontières~?
\end{interviewq}

\begin{interviewq}[$\star\star$ \iqroles{risque}]\label{iq:hf:latency-arbitrage-and-its-defence:5}
La liaison hertzienne de votre firme tombe pendant un orage. Qu'est-ce qui change dans le risque des stratégies, et que faites-vous~?
\end{interviewq}

\begin{interviewq}[$\star\star\star$ \iqroles{chercheur}]\label{iq:hf:latency-arbitrage-and-its-defence:6}
La latence d'annulation du fournisseur suit une loi exponentielle de moyenne $c$, et la latence de chacun des $n$ snipers une loi exponentielle de moyenne $s$, toutes ces variables étant indépendantes. Quelle est la
probabilité que le fournisseur soit victime de sniping~?
\end{interviewq}
